% Source: Corsin Nick/Class Notes/Week 12.pdf, pp. 5--7
\chapter{Week 12 --- Pullbacks, Integration of Forms \& Orientability}
\label{ch:week12}

% Transition: Claude Sonnet 5
This chapter extends last week's differential forms machinery in two directions needed to
state Stokes' theorem in general: the \newterm{pullback} $F^*\omega$ of a form under a smooth
map, which lets a form on the target be transported back to the source and is exactly the
operation underlying every change-of-variables computation so far, and the definition of
$\int_L \omega$ for a form on a genuinely curved (not just open-in-$\mathbb{R}^n$) submanifold
$L$. It closes with the notion of an \newterm{orientable manifold} --- an atlas whose
transition maps all preserve orientation --- quoted from the lecture script, which Corsin
annotates with a hand-drawn counterexample. The exercise sheet is at the end of the chapter,
in \Cref{sec:week12_solutions}.

\begin{ainote}
Corsin opens this week by recapping the exterior derivative $d : \Omega^k(U) \to
\Omega^{k+1}(U)$ --- its three defining properties, the notation $dx^i = e_i^*(x)$, and the
worked example computing $d\omega$ for the $1$-form $\omega_F$ associated to a vector field
$F : U \subseteq \mathbb{R}^3 \to \mathbb{R}^3$ --- verbatim from Week 11, pp.\ 13--15.
This exactly duplicates \cref{thm:exterior_derivative,ex:exterior_derivative_gradient_curl}
already transcribed there; not repeated here.
\end{ainote}

\section{Pullback of forms}
\label{sec:pullback_of_forms}

% Source: Corsin Nick/Class Notes/Week 12.pdf, p. 5
\begin{definition}[Pullback of a form]
\label{def:pullback_of_forms}
Let $U \subseteq \mathbb{R}^m$, $V \subseteq \mathbb{R}^n$ be open, $F : U \to V$ a smooth
map, and $\omega \in \Omega^p(V)$ a differential $p$-form on $V$. The \newterm{pullback}
(\germanterm{Zur\"uckziehung}) of $\omega$ under $F$ is the differential $p$-form on $U$
defined by
\[(F^*\omega)_x(v_1,\ldots,v_p) := \omega_{F(x)}\big(DF_x(v_1),\ldots,DF_x(v_p)\big).\]
\end{definition}

\begin{ainote}
Corsin also draws a commutative diagram $TU \xrightarrow{dF} TV \xrightarrow{\omega} \mathbb{R}$
over $U \xrightarrow{F} V$, with tangent-bundle projections $\pi_U, \pi_V$, and labels it
\qt{not exam relevant}. Omitted here for that reason; the formula above is the operative
content.
\end{ainote}

The pullback satisfies two properties that make it the right tool for transporting forms
along smooth maps:
\begin{enumerate}[label=\textbf{(\arabic*)}]
  \item $F^*(\omega\wedge\theta) = F^*\omega \wedge F^*\theta$;
  \item $F^*(d\omega) = d(F^*\omega)$ (the pullback commutes with the exterior derivative).
\end{enumerate}

\begin{ainote}
This is precisely the operation used, informally, in the reparametrization-invariance
motivation for \cref{def:antisymmetric_k_linear_form} in Week 11: writing $\phi^*\omega$ and
$(\phi\circ\Psi)^*\omega$ for the two sides of that calculation makes it read as $\int_D
\phi^*\omega = \int_V (\phi\circ\Psi)^*\omega = \int_V \Psi^*(\phi^*\omega)$, i.e.\ property (1)
applied to $\phi^*\omega$ and the fact that pullback along a composition is the composition of
pullbacks.
\end{ainote}

\section{Integration of differential forms on submanifolds}
\label{sec:integration_of_forms}

% Source: Corsin Nick/Class Notes/Week 12.pdf, p. 6
If $\mathcal{M}$ is an $m$-dimensional submanifold of $\mathbb{R}^m$ --- i.e.\ simply an open
set --- and $\omega \in \Omega^m(\mathcal{M})$ an $m$-form, then
\[\int_{\mathcal{M}} \omega := \int_{\mathcal{M}} \omega_x(e_1,\ldots,e_m)\,dx;\]
in particular, for $f \in C^\infty(\mathcal{M})$,
\[\int_{\mathcal{M}} f\,dx^1\wedge\cdots\wedge dx^m = \int_{\mathcal{M}} f(x)\,dx,\]
recovering the ordinary Lebesgue integral of $f$.

\begin{definition}[Integral of a form over a parametrized submanifold]
\label{def:integral_of_form_over_submanifold}
If $L^k \subseteq \mathbb{R}^n$ is a $k$-dimensional submanifold, $n>k$, and $\phi : U \to L$
is a parametrization with $U \subseteq \mathbb{R}^k$ open, then the integral of $\omega \in
\Omega^k(L)$ is defined as
\[\int_L \omega := \int_U \phi^*\omega = \int_U \omega_{\phi(x)}\!\left(\frac{\partial\phi}
  {\partial x^1},\ldots,\frac{\partial\phi}{\partial x^k}\right) dx.\]
\end{definition}

In particular, for $f \in C^\infty(L)$ and $1\leq i_1<\cdots<i_k\leq n$,
\[\int_L f\,dx^{i_1}\wedge\cdots\wedge dx^{i_k} = \int_U f(x)\det\!\left(\frac{\partial
  \phi^{i_s}}{\partial x^t}\right)_{1\leq s,t\leq k} dx,\]
because, by property (4) of \cref{def:wedge_product} in Week 11,
\[dx^{i_1}\wedge\cdots\wedge dx^{i_k}(\partial_1\phi,\ldots,\partial_k\phi) = \det\!\left(
  \begin{smallmatrix} dx^{i_1}(\partial_1\phi) & \cdots & dx^{i_1}(\partial_k\phi) \\
  \vdots & \ddots & \vdots \\ dx^{i_k}(\partial_1\phi) & \cdots & dx^{i_k}(\partial_k\phi)
  \end{smallmatrix}\right) = \det\!\left(\frac{\partial\phi^{i_s}}{\partial x^t}\right)_{1\leq
  s,t\leq k}.\]

\begin{ainote}
Applying \cref{def:integral_of_form_over_submanifold} to $\omega_F^{n-1}$ or $\omega_F^2$ from
Week 11's example recovers exactly the flux integrals of Week 10 --- those were, all along,
integrals of forms in this sense. And \cref{def:pullback_of_forms} is what makes the
motivating reparametrization-invariance calculation of Week 11 precise: $\int_L \omega$ does
not depend on the parametrization $\phi$ chosen, only on $L$ and its orientation.
\end{ainote}

\section{Orientable manifolds}
\label{sec:orientable_manifolds}

% Source: Corsin Nick/Class Notes/Week 12.pdf, p. 7
\begin{ainote}
The two definitions below are reproduced by Corsin directly from the lecture script (numbered
there as Definitions 14.68 and 14.65), rather than written out in his own hand; he adds the
margin annotations quoted and figured below.
\end{ainote}

\begin{definition}[Support of a $k$-form]
\label{def:support_of_k_form}
Let $U \subseteq \mathbb{R}^n$ be open and $\alpha \in \Omega^k(U)$. The \newterm{support}
(\germanterm{Tr\"ager}) of $\alpha$ is the closed set
\[\supp\alpha := \overline{\{x\in U : \alpha_x \neq 0\}} \subseteq \overline{U}.\]
We say $\alpha$ has \newterm{compact support} in $U$ if $\supp\alpha$ is bounded and contained
in $U$.
\end{definition}

\begin{definition}[Orientable manifold]
\label{def:orientable_manifold}
Let $M \subseteq \mathbb{R}^n$ be a $k$-dimensional $C^\infty$ submanifold. We say $M$ is
\newterm{orientable} (\germanterm{orientierbar}) if there exists a family of local
parametrizations $\phi_i : V_i \to M \cap U_i$, $i \in I$, whose images cover $M$, and such
that for every $i,j \in I$ with $M \cap U_i \cap U_j \neq \emptyset$, the transition map
$\widetilde\phi_i^{-1}\circ\widetilde\phi_j : \widetilde V_j \to \widetilde V_i$ has positive
Jacobian determinant:
\[\det D(\widetilde\phi_i^{-1}\circ\widetilde\phi_j) > 0 \qquad \text{on } \widetilde V_j.\]
Such a family $\{\phi_i\}$ is called an \newterm{atlas} of $M$; if it satisfies the
determinant condition above it is called an \newterm{orientation} of $M$. Once such a family
is chosen, $M$ is said to be \newterm{oriented}.
\end{definition}

\begin{ainote}
Corsin marks the family $\{\phi_i\}$ itself with the margin note \qt{is called an Atlas} next
to the covering condition, and the positivity condition with \qt{Atlas} again next to the
determinant inequality --- i.e.\ he is flagging that \qt{atlas} names the family regardless of
whether the determinant condition holds, while \qt{orientation} is the stronger notion once it
does. This matches the definition's own wording above.
\end{ainote}

Corsin illustrates the failure of the orientation condition with a hand-drawn counterexample:
two overlapping local parametrizations $\phi, \psi$ of the same patch of a surface, whose
coordinate frames $(\partial_1\phi,\partial_2\phi)$ and $(\partial_1\psi,\partial_2\psi)$ at a
shared point disagree in handedness.

% Sonnet 5 (Medium) --- FIG-W12-01
\begin{center}
\begin{tikzpicture}[>=Latex, scale=1.0]
  \draw[green!50!black, thick] (0,0) .. controls (1.5,0.9) and (3.5,0.6) .. (5,1.2)
    .. controls (5,2.4) and (3,2.6) .. (1.5,2.0) .. controls (0,1.6) and (0,0.8) .. (0,0);
  \draw[green!50!black, dotted] (0.4,1.0) .. controls (1.8,0.6) and (3.2,1.1) .. (4.4,1.6);
  \fill (2.1,1.3) circle (1.2pt);
  \draw[blue!70!black, thick, ->] (2.1,1.3) -- (3.3,1.55) node[right, font=\scriptsize]
    {$\partial_1\phi$};
  \draw[blue!70!black, thick, ->] (2.1,1.3) -- (1.7,2.15) node[above, font=\scriptsize]
    {$\partial_2\phi$};
  \draw[red, thick, ->] (2.1,1.3) -- (0.9,1.05) node[left, font=\scriptsize] {$\partial_1\psi$};
  \draw[red, thick, ->] (2.1,1.3) -- (2.55,0.35) node[below, font=\scriptsize] {$\partial_2\psi$};
  \node[font=\scriptsize] at (0.1,-0.35) {$\phi$};
  \node[font=\scriptsize] at (4.4,2.7) {$\psi$};
  \draw[blue!70!black, ->] (6.3,0.2) -- (7.3,0.2) node[right, font=\scriptsize] {$e_1$};
  \draw[blue!70!black, ->] (6.3,0.2) -- (6.3,1.2) node[above, font=\scriptsize] {$e_2$};
  \draw[red, ->] (8.6,0.2) -- (9.6,0.2) node[right, font=\scriptsize] {$e_1$};
  \draw[red, ->] (8.6,0.2) -- (8.6,1.2) node[above, font=\scriptsize] {$e_2$};
\end{tikzpicture}
\end{center}

\begin{ainote}
At the shared point, $(\partial_1\phi,\partial_2\phi)$ turns counterclockwise while
$(\partial_1\psi,\partial_2\psi)$ turns clockwise --- the two frames are mirror images of each
other. Corsin's caption: \qt{Here, $\det(D(\phi^{-1}\circ\psi)) < 0$!} This is exactly the atlas
$\{\phi,\psi\}$ failing \cref{def:orientable_manifold}: it covers the patch, but its transition
map has \emph{negative} Jacobian determinant there, so $\{\phi,\psi\}$ is an atlas but not an
orientation.
\end{ainote}

% Generator: Claude Sonnet 5 (Medium)
\begin{ainote}
\Cref{ex:12.3} and Exercise 12.6 (the latter not reproduced below, see the priority table)
invoke \newterm{Stokes' theorem} in its classical $\mathbb{R}^3$ form, which none of Corsin's
transcribed pages state explicitly. Recorded here for self-containedness, since the
exercise-sheet solutions below rely on it.
\end{ainote}

\begin{theorem}[Stokes' theorem, classical form in $\mathbb{R}^3$]
\label{thm:stokes_classical}
% Generator: Claude Sonnet 5 (Medium)
Let $S \subseteq \mathbb{R}^3$ be a compact, oriented $C^1$ surface with boundary $\partial S$
carrying the induced orientation, and let $F$ be a $C^1$ vector field on a neighborhood of
$S$. Then
\[\int_S \curl F \cdot d\vec{n} = \int_{\partial S} F \cdot d\vec{s}.\]
\end{theorem}

\exercisesheet{12}

\textit{Statements quoted verbatim from \texttt{exercises/Ex12\_Analysis2\_eng.pdf} (assigned
8 May 2026, due 18 May 2026). Attribution: Prof.\ Joaquim Serra, D-MATH, ETH Z\"urich.}

% Source: Corsin Nick/Class Notes/Week 12.pdf, p. 1
\begin{ainote}[Corsin's recommendations]
Colour key, as given on the page: blue \textbf{= important}, orange \textbf{= semi-important},
red \textbf{= optional}.

\begin{center}
\begin{tabular}{@{}l l@{}}
\textbf{Problem} & \textbf{Priority} \\
\hline
12.1 & semi-important \\
12.2 & optional \\
12.3 & important \\
12.4 & important \\
12.5 & important \\
12.6 & semi-important \\
12.7 & semi-important \\
12.8 & important \\
12.9 & important \\
\end{tabular}
\end{center}
\end{ainote}

\begin{ainote}
Only the exercises tagged important above are reproduced below: 12.3--12.5, 12.8, and 12.9.
12.1 (rigorous concatenation of paths), 12.2 (existence of a potential via a path integral,
building the tool used implicitly in \cref{ex:12.5} below), 12.6 (a Stokes reality check
structurally identical to \cref{ex:12.3} part 3), and 12.7 (four standard $\curl$/$\divg$
product-rule identities) are not reproduced, consistent with the standing decision to follow
Corsin's priority marking.
\end{ainote}

% Source: exercises/Ex12_Analysis2_eng.pdf, p. 1
\begin{exercise}[12.3 --- Work along paths and Stokes \textnormal{(important)}]
\label{ex:12.3}
Let $F : \mathbb{R}^3 \setminus (\{(0,0)\}\times\mathbb{R}) \to \mathbb{R}^3$ be the vector
field defined by
\[F(x,y,z) = \left(\frac{2(xz+y)}{x^2+y^2},\ \frac{2(yz-x)}{x^2+y^2},\ \log(x^2+y^2)\right)\]
and let $\gamma_1,\gamma_2,\gamma_3 : [0,2\pi] \to \mathbb{R}^3$ be the paths
\[\gamma_1(t) = (\cos t,\sin t,0), \qquad \gamma_2(t) = (2\cos t,2\sin t,2), \qquad
  \gamma_3(t) = (3\cos t+6,3\sin t,3).\]
\begin{enumerate}[label=\textbf{\arabic*.}]
  \item Describe the domain of definition of $F$. Is it connected? Is it simply connected?
  \item Compute the rotation $\curl(F)$ of $F$ and the path integral $\int_{\gamma_1} F\cdot
  d\gamma_1$ of $F$ along $\gamma_1$. Is $F$ conservative?
  \item Explain how the equation $\int_{\gamma_1} F\cdot d\gamma_1 = \int_{\gamma_2} F\cdot
  d\gamma_2$ follows from Stokes' theorem.
  \item Does Stokes' theorem imply that $\int_{\gamma_1} F\cdot d\gamma_1 = \int_{\gamma_3}
  F\cdot d\gamma_3$ holds?
\end{enumerate}
\end{exercise}

% Source: exercises/Ex12_Analysis2_eng.pdf, p. 2
\begin{exercise}[12.4 --- True or False \textnormal{(important)}]
\label{ex:12.4}
Let $U := \mathbb{R}^3\setminus\{0\}$ and $F \in C^1(U,\mathbb{R}^3)$. For $r>0$ set
\[I_r := \int_{\partial B_r} \langle F,n\rangle\,dS\]
where $n$ denotes the outward normal. Which of the following statements are always true?
\begin{enumerate}[label=\textbf{(\alph*)}]
  \item If $F$ is divergence-free, then $I_r=0$ for all $r>0$.
  \item If $F$ is divergence-free, then the value of $I_r$ does not depend on $r$.
  \item If $F$ is curl-free, then the value of $I_r$ does not depend on $r$.
  \item If $F = \curl(G)$ for a vector field $G \in C^1(U,\mathbb{R}^3)$ and $\lvert F(x)
  \rvert \leq C/\lvert x\rvert$, then $I_r=0$ for all $r>0$.
\end{enumerate}
\end{exercise}

% Source: exercises/Ex12_Analysis2_eng.pdf, p. 2
\begin{exercise}[12.5 --- Conservative vector fields \textnormal{(important)}]
\label{ex:12.5}
For which values of $\lambda \in \mathbb{R}$ is the vector field $F : \mathbb{R}^2 \to
\mathbb{R}^2$ defined by
\[F(x,y) = \begin{pmatrix} \lambda x e^y \\ (y+1+x^2)e^y \end{pmatrix}\]
conservative? Determine a potential for $F$ for these values.
\end{exercise}

% Source: exercises/Ex12_Analysis2_eng.pdf, p. 3
\begin{exercise}[12.8 --- Multiple Choice \textnormal{(important)}]
\label{ex:12.8}
Let $Q = [a,b]\times[c,d] \subseteq \mathbb{R}^2$ and $f \in C(Q,\mathbb{R}^2)$. Let $n$ be an
outward normal field on $Q$ and
\[I_b = \int_{[a,b]\times\{c\}} \langle f,n\rangle, \qquad I_t = \int_{[a,b]\times\{d\}}
  \langle f,n\rangle, \qquad I_l = \int_{\{a\}\times[c,d]} \langle f,n\rangle, \qquad
  I_r = \int_{\{b\}\times[c,d]} \langle f,n\rangle.\]
Which of the following expressions corresponds to the outer flux of $f$ through $\partial Q$?
\begin{enumerate}[label=\textbf{(\alph*)}]
  \item $I_b - I_t + I_r - I_l$
  \item $I_b + I_t + I_r + I_l$
  \item $I_b + I_t - I_r - I_l$
  \item $-I_b + I_t - I_r + I_l$
\end{enumerate}
\end{exercise}

% Source: exercises/Ex12_Analysis2_eng.pdf, p. 3
\begin{exercise}[12.9 --- Change of variables with differential forms \textnormal{(important)}]
\label{ex:12.9}
Let $U,V \subseteq \mathbb{R}^n$ be open sets, and let $\psi = (\psi^1,\ldots,\psi^n) : U \to
V$ be a diffeomorphism with $\det\!\left(\frac{\partial\psi^j}{\partial x^i}\right) > 0$.
Define an $n$-differential form $\omega$ on $U$ by $\omega = d\psi^1\wedge d\psi^2\wedge\cdots
\wedge d\psi^n$.
\begin{enumerate}[label=\textbf{(\alph*)}]
  \item Show that $\omega = \det\!\left(\frac{\partial\psi^j}{\partial x^i}\right) dx^1\wedge
  \cdots\wedge dx^n$.
  \item From the relation in (a), rewrite the change of variables formula using differential
  forms.
\end{enumerate}
\end{exercise}

\section{Solutions}
\label{sec:week12_solutions}

\begin{exercisesolution}[Solution of \cref{ex:12.3}]
\textbf{1.} The domain is $\{(x,y,z) \in \mathbb{R}^3 : (x,y) \neq (0,0)\}$, i.e.\
$\mathbb{R}^3$ with the $z$-axis removed. It is connected (path-connected: any two points can
be joined by a path avoiding the axis), but not simply connected --- e.g.\ $\gamma_1$, which
winds once around the axis, cannot be contracted to a point without crossing it. (It is also
not convex: the segment from $(-1,0,0)$ to $(1,0,0)$ leaves the domain.)

\textbf{2.} Writing $F = (F^1,F^2,F^3)$ with $F^1 = 2(xz+y)/(x^2+y^2)$, $F^2 =
2(yz-x)/(x^2+y^2)$, $F^3 = \log(x^2+y^2)$, a direct computation of each pair of mixed
partials gives
\[\partial_y F^3 = \partial_z F^2 = \frac{2y}{x^2+y^2}, \qquad
  \partial_z F^1 = \partial_x F^3 = \frac{2x}{x^2+y^2}, \qquad
  \partial_x F^2 = \partial_y F^1 = \frac{2(x^2-y^2-2xyz)}{(x^2+y^2)^2},\]
so $\curl F = (\partial_yF^3-\partial_zF^2,\ \partial_zF^1-\partial_xF^3,\ \partial_xF^2-
\partial_yF^1) = (0,0,0)$: $F$ is curl-free. On $\gamma_1$, $x=\cos t$, $y=\sin t$, $z=0$,
$x^2+y^2=1$, so $F(\gamma_1(t)) = (2\sin t,\,-2\cos t,\,0)$ and $\gamma_1'(t) = (-\sin t,
\cos t,0)$, giving
\[\int_{\gamma_1} F\cdot d\gamma_1 = \int_0^{2\pi} \big(2\sin t\cdot(-\sin t) + (-2\cos t)
  \cdot\cos t\big)\,dt = \int_0^{2\pi} -2\,dt = -4\pi.\]
Since this is nonzero on a closed loop, $F$ is \textbf{not} conservative --- even though it is
curl-free, its domain is not simply connected (part 1), so curl-free does not force
conservative here.

\textbf{3.} Let $S := \{(x,y,z) : \sqrt{x^2+y^2} = 1+z/2,\ 0\leq z\leq2\}$, the cone swept out
by linearly interpolating the radius from $\gamma_1$ ($z=0$, radius $1$) to $\gamma_2$ ($z=2$,
radius $2$), parametrized by $H(s,t) := \big((1+s)\cos t,(1+s)\sin t,2s\big)$ for $(s,t) \in
[0,1]\times[0,2\pi]$. Since the radius $1+s \geq 1 > 0$ throughout, $S$ never meets the
$z$-axis and lies entirely in the domain of $F$; its oriented boundary (outward normal) is
$\gamma_2$ minus $\gamma_1$. Since $\curl F = 0$ on $S$ by part 2, \cref{thm:stokes_classical}
gives
\[\int_{\gamma_2} F\cdot d\gamma_2 - \int_{\gamma_1} F\cdot d\gamma_1 = \int_{\partial S}
  F\cdot ds = \int_S \curl F\cdot d\vec n = 0,\]
i.e.\ $\int_{\gamma_1} F\cdot d\gamma_1 = \int_{\gamma_2} F\cdot d\gamma_2$.

\textbf{4.} No. The disk $D := \{(x,y,z) : (x-6)^2+y^2 \leq 9,\ z=3\}$, oriented upward, has
oriented boundary exactly $\gamma_3$, and $D$ lies entirely in the domain of $F$ (it does not
meet the $z$-axis, since its center $(6,0,3)$ is at distance $6$ from the axis and its radius
is only $3$). By \cref{thm:stokes_classical} with $\curl F=0$,
\[\int_{\gamma_3} F\cdot d\gamma_3 = \int_{\partial D} F\cdot ds = \int_D \curl F\cdot d\vec n
  = 0.\]
Since $\int_{\gamma_1} F\cdot d\gamma_1 = -4\pi \neq 0 = \int_{\gamma_3} F\cdot d\gamma_3$, the
two do \emph{not} agree, and Stokes' theorem gives no reason to expect they would: unlike the
pair $(\gamma_1,\gamma_2)$, there is no surface joining $\gamma_1$ and $\gamma_3$ that stays
inside the domain, since $\gamma_1$ winds around the excluded axis and $\gamma_3$ does not.
\end{exercisesolution}

\begin{ainote}
Solved independently before consulting \texttt{Sol12\_Analysis2\_eng.pdf}; all four parts,
including the explicit cone parametrization in part 3 and the disk in part 4, match the
official solution exactly (the official solution's part 1 also remarks on non-convexity, added
above for the same reason). One phrasing note: the official solution's part 2 labels the
vanishing of $\curl F$ as \qt{$F$ is divergence-free} where it visibly means curl-free
(computes and displays $\curl F$, not $\divg F$); read as a wording slip, not a substantive
error, and not repeated here.
\end{ainote}

\begin{exercisesolution}[Solution of \cref{ex:12.4}]
\textbf{(a) False.} The Newton potential's gradient $F(x) = \nabla(1/\lvert x\rvert) =
-x\lvert x\rvert^{-3}$ is divergence-free on $U$ (it is $V_{-3}$ up to sign, from
\cref{ex:11.5} of Week 11, with $n=3$), yet $I_r = \int_{\partial B_r} F\cdot n\,dS =
-\vol_2(\partial B_r) \neq 0$.

\textbf{(b) True.} Applying \cref{thm:gauss_divergence} on the shell $B_{r_2}\setminus
B_{r_1} \Subset U$ for $0<r_1<r_2$, whose boundary is $\partial B_{r_2}$ (outward) together
with $\partial B_{r_1}$ with reversed orientation,
\[0 = \int_{B_{r_2}\setminus B_{r_1}} \divg F\,dx = I_{r_2} - I_{r_1},\]
so $I_r$ is independent of $r$.

\textbf{(c) False.} Take $F(x) = x = \tfrac12\nabla\lvert x\rvert^2$, which is curl-free
(being a gradient). Then $\divg F \equiv 3$, so by \cref{thm:gauss_divergence}, $I_r =
\int_{B_r} 3\,dx = 3\vol_3(B_r) = 4\pi r^3$, which depends on $r$.

\textbf{(d) True.} If $F = \curl G$ then $\divg F = \divg(\curl G) = 0$ (a standard vector
identity), so part (b) applies and $I_r$ is independent of $r$. On the other hand,
\[\lvert I_r\rvert \leq \int_{\partial B_r} \lvert\langle F,n\rangle\rvert\,dS \leq
  \frac{C}{r}\,\vol_2(\partial B_r) = \frac{C}{r}\cdot4\pi r^2 = 4\pi C r \longrightarrow 0
  \quad\text{as } r\to0^+.\]
Since $I_r$ is constant in $r$ and this constant is a limit of values tending to $0$, $I_r=0$
for all $r>0$.
\end{exercisesolution}

\begin{ainote}
Solved independently before consulting \texttt{Sol12\_Analysis2\_eng.pdf}; all four parts
match, including the same Newton-potential counterexample for (a) and the same $F=x$
counterexample for (c).
\end{ainote}

\begin{exercisesolution}[Solution of \cref{ex:12.5}]
Writing $F = (F^1,F^2)$ with $F^1 = \lambda xe^y$, $F^2 = (y+1+x^2)e^y$, the integrability
condition is $\partial_yF^1 = \partial_xF^2$:
\[\partial_yF^1 = \lambda xe^y, \qquad \partial_xF^2 = 2xe^y.\]
Since $\mathbb{R}^2$ is convex (so integrability is equivalent to conservativity here), $F$
is conservative if and only if $\lambda=2$.

For $\lambda=2$, seek $f$ with $\partial_xf = F^1 = 2xe^y$: integrating in $x$,
$f(x,y) = x^2e^y + g(y)$ for some function $g$. Then $\partial_yf = x^2e^y+g'(y)$ must equal
$F^2 = (y+1+x^2)e^y$, so $g'(y) = (y+1)e^y$; integrating by parts, $\int(y+1)e^y\,dy =
ye^y - e^y + e^y = ye^y$ (up to an additive constant). Hence
\[f(x,y) = \boxed{\ (x^2+y)e^y\ } \qquad (\text{unique up to an additive constant}),\]
which one checks directly satisfies $\nabla f = F$.
\end{exercisesolution}

\begin{ainote}
Solved independently before consulting \texttt{Sol12\_Analysis2\_eng.pdf}; both the value
$\lambda=2$ and the potential $(x^2+y)e^y$ match the official solution exactly.
\end{ainote}

\begin{exercisesolution}[Solution of \cref{ex:12.8}]
The four quantities $I_b,I_t,I_l,I_r$ are already defined using the same outward normal field
$n$ on each of the four edges of $\partial Q$, so the total outer flux is simply their sum,
with no relative sign corrections needed: $\boxed{I_b+I_t+I_r+I_l}$, answer \textbf{(b)}.
\end{exercisesolution}

\begin{ainote}
Solved independently before consulting \texttt{Sol12\_Analysis2\_eng.pdf}; matches (answer
(b), no further justification given there either).
\end{ainote}

\begin{exercisesolution}[Solution of \cref{ex:12.9}]
\textbf{(a)} Writing each $d\psi^j$ in the coordinate basis, $d\psi^j = \sum_{i=1}^n
\frac{\partial\psi^j}{\partial x^i}\,dx^i$, so
\[\omega = \left(\sum_{i_1=1}^n \frac{\partial\psi^1}{\partial x^{i_1}}dx^{i_1}\right)\wedge
  \cdots\wedge\left(\sum_{i_n=1}^n \frac{\partial\psi^n}{\partial x^{i_n}}dx^{i_n}\right)
  = \sum_{i_1,\ldots,i_n} \frac{\partial\psi^1}{\partial x^{i_1}}\cdots
  \frac{\partial\psi^n}{\partial x^{i_n}}\,dx^{i_1}\wedge\cdots\wedge dx^{i_n}\]
by bilinearity of $\wedge$ (property (1) of \cref{def:wedge_product}). Any term with a
repeated index vanishes, since $dx^i\wedge dx^i=0$ by antisymmetry, so only tuples
$(i_1,\ldots,i_n) = (\sigma(1),\ldots,\sigma(n))$ for a permutation $\sigma \in S_n$
contribute; using $dx^{\sigma(1)}\wedge\cdots\wedge dx^{\sigma(n)} = \sgn(\sigma)\,dx^1\wedge
\cdots\wedge dx^n$,
\[\omega = \left(\sum_{\sigma\in S_n} \sgn(\sigma)\,\frac{\partial\psi^1}{\partial
  x^{\sigma(1)}}\cdots\frac{\partial\psi^n}{\partial x^{\sigma(n)}}\right) dx^1\wedge\cdots
  \wedge dx^n.\]
The coefficient is exactly the Leibniz expansion of $\det(D\psi)$, so
\[\boxed{\ d\psi^1\wedge\cdots\wedge d\psi^n = \det(D\psi)\,dx^1\wedge\cdots\wedge dx^n.\ }\]

\textbf{(b)} Let $\eta = f(y)\,dy^1\wedge\cdots\wedge dy^n$ be an $n$-form on $V$. Its
pullback under $\psi$ is $\psi^*\eta = f(\psi(x))\,d\psi^1\wedge\cdots\wedge d\psi^n$, which by
part (a) equals $f(\psi(x))\det(D\psi(x))\,dx^1\wedge\cdots\wedge dx^n$. By
\cref{def:integral_of_form_over_submanifold} (applied with $U$ itself as the $n$-dimensional
submanifold),
\[\int_U \psi^*\eta = \int_U f(\psi(x))\det(D\psi(x))\,dx.\]
Since $\det(D\psi) = \lvert\det(D\psi)\rvert>0$ by hypothesis, the right-hand side is exactly
the ordinary change-of-variables formula, so
\[\int_U \psi^*\eta = \int_U f(\psi(x))\,\lvert\det(D\psi(x))\rvert\,dx = \int_V f(y)\,dy
  = \int_V \eta.\]
Thus $\boxed{\int_U \psi^*\eta = \int_V \eta}$: the change-of-variables formula is exactly the
statement that integration of forms is invariant under pullback by an orientation-preserving
diffeomorphism.
\end{exercisesolution}

\begin{ainote}
Solved independently before consulting \texttt{Sol12\_Analysis2\_eng.pdf}; both parts match
the official solution, including the identification of (b) with pullback-invariance of form
integration.
\end{ainote}
